\chapter{Cat and Dog Allergy}
\index{Cat and Dog Allergy}

\section*{Definition and Overview}
Cat and dog allergy is an allergic reaction to proteins in the skin flakes (dander), saliva, and urine of pet cats and dogs. When a person with the allergy breathes in or touches these proteins, the immune system overreacts and produces symptoms affecting the nose, eyes, airways, and skin.

The main allergen in cats is a protein called Fel d 1, found in the skin and saliva and spread through grooming and shedding. Dog allergens include Can f 1 and other proteins found in dander, saliva, and urine. These tiny, sticky particles float in the air, settle on furniture, carpets, and clothing, and can remain in a home long after the pet has gone. Pet allergy is common in people with other allergies and can worsen asthma.

\section*{Epidemiology}
Allergy to cats and dogs is one of the most common animal allergies. Roughly one in five to one in three people with allergic disorders react to pets, and about 10--20\% of the general population is sensitised to cat or dog allergens even if they do not notice symptoms. Because many households keep pets, exposure is almost universal.

The allergy can appear at any age, though it often develops in childhood or early adulthood and may run in families. People with asthma, hay fever, or eczema are more likely to have pet allergy. Some research suggests that children raised with pets from infancy may be less likely to develop allergies, but the evidence is not clear-cut, and in a person who is already allergic, pets make symptoms worse.

\section*{Aetiology and Pathophysiology}
Pet allergy is a type I hypersensitivity reaction, meaning it is driven by immunoglobulin E (IgE) antibodies. The steps are:

\begin{itemize}
    \item \textbf{Sensitisation:} the immune system of a susceptible person mistakes a pet protein, such as Fel d 1 or Can f 1, for a threat and produces specific IgE antibodies against it
    \item \textbf{Re-exposure:} the next time the person meets the allergen, it binds to these antibodies on cells called mast cells
    \item \textbf{Chemical release:} mast cells release histamine and other inflammatory chemicals
    \item \textbf{Symptoms:} these chemicals cause a runny nose, sneezing, itchy eyes, cough, and narrowing of the airways in people with asthma
\end{itemize}

Pet allergens are carried on small particles that stay airborne for long periods and cling to soft furnishings, so symptoms can continue even when the pet is not present. Cat allergen in particular is very small and sticky and travels widely, including into homes without cats.

\section*{Clinical Features}
\begin{itemize}
    \item Sneezing, a runny or blocked nose, and an itchy nose
    \item Itchy, red, or watery eyes
    \item A scratchy or sore throat, cough, and a hoarse voice
    \item Wheezing, chest tightness, or shortness of breath -- especially in people with asthma
    \item Worsening asthma control, with more frequent use of reliever medication
    \item Hives or an itchy rash where a cat or dog has scratched or licked the skin
    \item Swelling and itching of the skin after contact with saliva or dander
    \item Symptoms triggered or worse in homes with pets, after visiting pet owners, or near bedding, carpets, and upholstery
    \item Sinus congestion, headache, or disturbed sleep
\end{itemize}

Symptoms may appear within minutes of exposure or develop over a few hours.

\section*{Differential Diagnosis}
\begin{itemize}
    \item \textbf{Hay fever (seasonal allergic rhinitis):} allergy to grass or tree pollen, often seasonal
    \item \textbf{Dust mite allergy:} similar symptoms, triggered by house dust mites in bedding and carpets
    \item \textbf{Mould or pollen allergy} that coexists with pet allergy
    \item \textbf{Non-allergic rhinitis:} nasal symptoms from irritants, temperature changes, or medicines
    \item \textbf{Respiratory infections:} colds, sinusitis, or viral infections causing similar nasal symptoms
    \item \textbf{Asthma and other triggers:} exercise, smoke, or cold air can cause wheeze unrelated to pets
    \item \textbf{Skin conditions:} eczema or contact dermatitis that might be mistaken for allergy to pet contact
\end{itemize}

Because several allergies and irritants commonly overlap, identifying pet allergy usually requires a careful history and allergy testing.

\section*{When to Seek Medical Care}
See a doctor if symptoms are persistent, distressing, or interfering with sleep, school, work, or daily life, or if you think you might benefit from allergy testing. People with asthma should have their asthma management reviewed if symptoms are worsening.

\textbf{Contact your local emergency services for:}
\begin{itemize}
    \item Severe difficulty breathing or wheezing that does not improve with a reliever puffer
    \item Chest tightness, blue lips, or struggling to speak
    \item Signs of a severe allergic reaction (anaphylaxis): swelling of the lips, tongue, or throat, or collapse
\end{itemize}

Anaphylaxis to pet allergens is rare but possible, particularly after a scratch or bite.

\section*{Diagnosis and Investigations}
\begin{itemize}
    \item \textbf{History:} a detailed account of symptoms and their link to pet exposure is the most important step
    \item \textbf{Allergy testing:} skin prick tests, in which tiny amounts of cat or dog allergen are placed on the skin, usually confirm the diagnosis quickly
    \item \textbf{Blood tests:} specific IgE (formerly RAST) blood tests can measure antibodies to cat and dog allergens
    \item \textbf{Lung function tests:} spirometry if asthma is suspected, including before and after reliever use
    \item \textbf{Trial of avoidance:} symptoms that clearly improve when away from pets support the diagnosis
    \item \textbf{Referral:} an allergy specialist may be involved if tests are unclear or symptoms are severe
\end{itemize}

\section*{Management}
\begin{itemize}
    \item \textbf{Avoidance:} reducing exposure is the most effective measure; where a pet must stay, create pet-free zones, especially the bedroom
    \item \textbf{Environmental control:} HEPA air filters, vacuuming with a HEPA filter, washing pet bedding, and removing or covering soft furnishings and carpets
    \item \textbf{Pet hygiene:} regular bathing of the pet and washing hands and clothes after contact, though these help only partly
    \item \textbf{Antihistamines:} non-drowsy antihistamine tablets or syrups relieve sneezing, itching, and a runny nose
    \item \textbf{Nasal sprays:} steroid nasal sprays reduce nasal inflammation and congestion with regular use
    \item \textbf{Eye drops:} antihistamine or anti-inflammatory drops for itchy, red eyes
    \item \textbf{Asthma medicines:} reliever and preventer inhalers, reviewed by a doctor, for people with asthma
    \item \textbf{Allergen immunotherapy:} allergy shots or under-the-tongue treatment can gradually reduce sensitivity and are considered for people with severe, poorly controlled symptoms
\end{itemize}

\section*{Complications}
\begin{itemize}
    \item Poorly controlled asthma, with more frequent flare-ups and a higher risk of asthma attacks
    \item Chronic nasal congestion leading to sinusitis or ear infections
    \item Disturbed sleep and daytime tiredness from night-time symptoms
    \item Reduced quality of life, including avoidance of friends and family who keep pets
    \item Hives and skin irritation from scratches and licks
    \item The emotional difficulty of giving up a beloved pet, which for many families is a major decision
\end{itemize}

\section*{Prognosis and Outcomes}
Pet allergy is usually a long-term condition that does not completely resolve, but symptoms are very manageable in most people. With good avoidance and the right medicines, most people achieve good control of their symptoms. In people with asthma, controlling pet allergy significantly improves asthma control. Allergen immunotherapy offers the possibility of reduced sensitivity over several years of treatment. Some people find their symptoms change over time, improving in some seasons or worsening with other allergies.

\section*{Prevention and Self-Care}
\begin{itemize}
    \item If you are allergic and have not yet got a pet, consider a pet that is less likely to trigger your symptoms, or avoid pets altogether
    \item Keep pets out of the bedroom and off beds and sofas to create an allergen-free sleeping area
    \item Use HEPA air purifiers and vacuum regularly with a HEPA-filter vacuum cleaner
    \item Wash the pet weekly if possible, and wash your hands after touching the animal
    \item Remove carpets and heavy curtains from bedrooms, and wash bedding in hot water
    \item Take allergy medicines before visiting homes with pets, and shower and change clothes afterwards
    \item Keep asthma well controlled with regular preventer use and a written asthma action plan
    \item Have allergy testing and advice if you are unsure what is causing your symptoms
\end{itemize}

\section*{Sources and Further Reading}
The following organisations provide reliable information on pet allergies and allergic disease.

\begin{itemize}
    \item National Health Service. \textit{NHS conditions guide on allergies and animal allergy.} https://www.nhs.uk/conditions/
    \item Mayo Clinic. \textit{Pet allergy: symptoms, causes, and treatment.} https://www.mayoclinic.org/
    \item Centers for Disease Control and Prevention. \textit{Information on pet allergens and environmental health.} https://www.cdc.gov/
    \item World Health Organization. \textit{Global guidance on allergic diseases and asthma.} https://www.who.int/
    \item MedlinePlus. \textit{Consumer health information on allergies and allergens.} https://medlineplus.gov/
    \item MSD Manual. \textit{Allergy overview and mechanisms -- professional reference.} https://www.msdmanuals.com/
\end{itemize}
